\section{Implications}
\label{sec:discussion}

\paragraph{Structural Linkability and Its Limits (\texorpdfstring{\hyperref[sec:question1]{RQ1}}{RQ1})}

The results demonstrate that interoperability of regulatory substance data is fundamentally constrained by the absence of a shared and formalised notion of chemical identity. While structure-based identifiers such as InChIKey enable unambiguous linking across datasets, a substantial fraction of regulatory entries cannot be expressed in this form. Instead, regulatory “substances” frequently represent heterogeneous entities, including mixtures, substance groups, and analytical parameters. This reflects a structural mismatch between chemical reality and regulatory abstraction: without explicit formalisation of equivalence, subset, and overlap relations, automated integration remains inherently limited.

\paragraph{Unreliability of CAS-based Identification (\texorpdfstring{\hyperref[sec:question2]{RQ2}}{RQ2})}

The analysis confirms that commonly used identifiers in regulatory contexts, particularly CAS numbers, do not provide a reliable foundation for interoperability. Although widely adopted, they exhibit both one-to-many and many-to-one mappings with chemical structures, reflecting differences in stereochemistry, salt forms, and context-dependent definitions. This finding challenges the widespread assumption in regulatory data integration that such identifiers can serve as globally unique references. In practice, relying on them introduces ambiguity rather than resolving it, undermining their suitability as a primary reference mechanism for cross-dataset linking.

\paragraph{Failure of Semantic Compensation (\texorpdfstring{\hyperref[sec:question3]{RQ3}}{RQ3})}

Embedding-based methods and large language models cannot compensate for the absence of structure-based identifiers. The limitations are inherent to the data: textual descriptions of mixtures, groups, and analytical entries lack the information required to infer molecular structure, so semantic similarity reflects lexical overlap rather than chemical identity. Advanced tooling raises the ceiling but does not dissolve this constraint.

\paragraph{The Combinatorial Barrier to Retrospective Harmonisation}

The scale of the interoperability problem becomes evident when considering the effort required for manual reconciliation. For non-structure substances, establishing semantic relationships requires pairwise comparison, leading to quadratic growth in workload. The estimated effort of several thousand person-years demonstrates that retrospective correction is not a viable strategy at scale, interoperability must be built into regulatory data models by design.

\paragraph{Towards Ontology-Driven Regulatory Data Infrastructures}

The 304 high-confidence SKOS triples illustrate concretely what ontology alignment enables: via \texttt{skos:broadMatch} ($n = 285$), a regulatory description such as \textit{fatty acids, C16--22} acquires machine-interpretable scope and its candidate member substances become enumerable through the ChemOnt hierarchy; via \texttt{skos:exactMatch} ($n = 18$), the ChemOnt class directly and completely defines the regulated group. For structure-defined substances, InChIKey-based linking and ChEBI annotations further enable cross-domain queries and composite scoring that are infeasible on fragmented datasets. However, the crosswalk covers only entries whose names contain sufficient chemical information, highlighting that the interoperability benefit of ontologies is maximised when group boundaries are defined explicitly in the regulatory text itself.

\paragraph{Implications for FAIR Chemical Data in Europe}

The findings have direct implications for the implementation of FAIR~(Findable, Accessible, Interoperable, Reusable) principles in European chemical regulation. Current regulatory datasets fall short of interoperability primarily because they lack globally unique, persistent, and machine-readable identifiers grounded in chemical structure. Achieving FAIR chemical substance data requires a shift from document-centric and identifier-fragmented approaches towards linked data infrastructures~\cite{BernersLee2006LinkedData,Heath2011LinkedData} in which substances are represented as resolvable URIs, grounded in structure-based identifiers and connected through formal ontologies. This shift is not purely technical: it requires alignment between data modelling practices and the legal definitions of substances in regulatory frameworks.

\paragraph{Limitations and Scope}

This study focuses on datasets derived from ECHA and related European regulatory sources, and on the use of InChIKey as the primary structure-based identifier. While the results are consistent across the analysed datasets, further work is needed to assess the generality of these findings across other regulatory domains, jurisdictions, and substance categories. Furthermore, not all regulatory concepts can or should be reduced to individual molecular structures. Mixtures, UVCBs, and analytical parameters are legitimate and necessary constructs in regulation. However, their current representation lacks formal semantics, limiting their integration with structure-based systems. Future work should explore hybrid models that explicitly represent these entities as sets, classes, or functions over well-defined chemical structures.

\paragraph{From Data Integration to Regulatory Design}

The interoperability problem is not primarily a technical integration challenge, but a consequence of how substances are defined and represented in regulation. Efforts to improve interoperability must therefore extend beyond data engineering to include the principled design of regulatory concepts, vocabularies, and data standards. Embedding structure-based identifiers, ontology references, and machine-readable semantics directly into regulatory processes~\cite{EIF2017,InteroperableEuropeAct2024} would significantly reduce ambiguity and enable scalable integration. Without such changes, the gap between regulatory data and scientific data infrastructures will continue to widen, limiting the effectiveness of data-driven environmental governance. In practice, this implies that regulatory data infrastructures should:
(i) include structure-based identifiers (e.g.~InChIKey) wherever applicable;
(ii) represent substance groups and mixtures using explicit semantic models (e.g.~set- or class-based representations);
(iii) publish substance data as linked data with persistent URIs.